\section{Overview}

This blueprint explains the Lean 4 companion project for the paper
\emph{A Sharper Hoeffding Bound for Weighted Sums of Exchangeable Random
Variables}.  The paper proves a Hoeffding-type moment generating function bound
for bounded exchangeable random variables:
\[
  \E \exp\left\{\lambda \sum_{i=1}^{n}
  w_i(X_i-\overline X_N)\right\}
  \le
  \exp\left\{
    \frac{\lambda^2}{2}\Gamma_N
    \|P_{\mathbf 1^\perp}\widetilde w\|_2^2
  \right\}.
\]
The finite-population inflation factor satisfies
\[
  \Gamma_N = 1+\frac{3}{2N}+O(N^{-2}).
\]

The Lean code is now modularized under
\[
  \texttt{Testlean/Exchangeable\_Hoeffding\_ver2}.
\]
The import
\[
  \texttt{import Testlean.Exchangeable\_Hoeffding\_ver2.All}
\]
loads the full development.  The declarations live in the namespace
\texttt{ExchangeableHoeffding}.

The project does not use global axioms.  Instead, the files record the paper's
major mathematical ingredients as explicit theorem and lemma interfaces.  Some
finite algebra, such as the zero-sum property of the projection, is already
proved directly in Lean.  The deeper analytic steps are stated as fields of
small structures:
\[
\begin{gathered}
  \texttt{MainProofObligations},\quad
  \texttt{SliceProofObligations},\quad
  \texttt{HypergeometricProofObligations},\\
  \texttt{ThreePointProofObligations},\quad
  \texttt{GammaProofObligations}.
\end{gathered}
\]
This makes the boundary between checked reductions and future analytic
formalization explicit.

\section{Module Map}

\begin{definition}[Full import]
\lean{ExchangeableHoeffding.theorem_2_1}
The module
\[
  \texttt{Testlean.Exchangeable\_Hoeffding\_ver2.All}
\]
imports the six working modules:
\[
\begin{array}{ll}
\texttt{Basic} & \text{finite vectors, projection, exchangeability},\\
\texttt{Gamma} & \Gamma_N,\ \text{closed form, Barber comparison},\\
\texttt{Hypergeometric} & \text{hypergeometric PMF/MGF and } B_{N,m},\\
\texttt{ThreePoint} & \text{three-coordinate extremizer objects},\\
\texttt{HammingSlice} & \text{slices, elementary symmetric polynomials},\\
\texttt{Main} & \text{main theorem, tail bound, optimality interface}.
\end{array}
\]
\end{definition}

\section{Formula--Lean Bridge}

The following table is the shortest way to compare the paper notation with the
Lean objects.

\[
\begin{array}{lll}
\text{Paper notation} & \text{Meaning} & \text{Lean declaration}\\
\hline
x\in\R^N & \text{finite vector} & \texttt{Fin N -> Real}\\
\widetilde w & \text{zero-padded weight vector} & \texttt{zeroPad hn w}\\
\overline X_N & \text{finite-population average} & \texttt{avg x}\\
P_{\mathbf 1^\perp}z & \text{projection off constants} & \texttt{projOnePerp z}\\
\sum_i w_i(X_i-\overline X_N) & \text{centered statistic} &
  \texttt{centeredWeightedSum hn w x}\\
\Gamma_N & \text{inflation factor} & \texttt{Gamma N}\\
F_{N,k}(y) & \text{Hamming-slice log-MGF} &
  \texttt{sliceFunctional k y}\\
B_{N,m} & \text{hypergeometric factor} & \texttt{hypergeomB N m}.
\end{array}
\]

\begin{definition}[Finite vectors and centering]
\lean{ExchangeableHoeffding.dot, ExchangeableHoeffding.normSq,
ExchangeableHoeffding.avg, ExchangeableHoeffding.projOnePerp,
ExchangeableHoeffding.zeroPad, ExchangeableHoeffding.centeredWeightedSum}
In \texttt{Basic.lean}, vectors are represented as functions
\(\operatorname{Fin} N\to\R\).  The dot product and squared norm are finite
sums:
\[
  \langle a,x\rangle=\sum_i a_i x_i,
  \qquad
  \|a\|_2^2=\sum_i a_i^2.
\]
The projection used in the paper is
\[
  P_{\mathbf 1^\perp}z
  =
  z-\frac{1}{N}\left(\sum_j z_j\right)\mathbf 1,
\]
which is the Lean declaration \texttt{projOnePerp}.  The centered statistic is
defined in Lean by the projected dot product
\[
  \texttt{centeredWeightedSum hn w x}
  =
  \left\langle P_{\mathbf 1^\perp}\widetilde w,x\right\rangle.
\]
\end{definition}

\begin{lemma}[Projection has zero sum]
\label{lem:projection-zero-sum}
\lean{ExchangeableHoeffding.sum_projOnePerp_eq_zero}
If \(N>0\), then
\[
  \sum_i (P_{\mathbf 1^\perp}z)_i=0.
\]
This finite-sum identity is already proved directly in Lean.
\end{lemma}

\begin{definition}[Exchangeable law]
\lean{ExchangeableHoeffding.permute, ExchangeableHoeffding.Exchangeable,
ExchangeableHoeffding.exchangeability_supports_symmetrization,
ExchangeableHoeffding.BoundedByOne}
The paper assumption \(X_\pi\stackrel d=X\) is represented by equality of
pushforward measures:
\[
  \texttt{Measure.map (permute sigma) mu = mu}.
\]
The theorem \texttt{exchangeability\_supports\_symmetrization} is the Lean name
for the direct use of this assumption.
\end{definition}

\section{Main Theorem}

\begin{definition}[Inflation factor]
\label{def:gamma}
\lean{ExchangeableHoeffding.gammaTerm, ExchangeableHoeffding.Gamma,
ExchangeableHoeffding.GammaClosed, ExchangeableHoeffding.harmonic,
ExchangeableHoeffding.epsilonBarber}
The paper's finite-population factor is
\[
  \Gamma_N
  =
  \max_{1\le s\le \lfloor N/2\rfloor}
  \frac{N(N-s)}{s}
  \sum_{\ell=N-s}^{N-1}\frac{1}{\ell^2}.
\]
The file \texttt{Gamma.lean} also defines the closed-form expression
\texttt{GammaClosed N}, the harmonic number, and Barber's comparison factor
\(\varepsilon_N\).
\end{definition}

\begin{lemma}[Endpoint comparison]
\lean{ExchangeableHoeffding.gamma_two, ExchangeableHoeffding.harmonic_two,
ExchangeableHoeffding.epsilonBarber_two,
ExchangeableHoeffding.gamma_eq_barber_two}
Lean checks the endpoint computations
\[
  \Gamma_2=2,\qquad H_2=\frac32,\qquad \varepsilon_2=1,
\]
and therefore
\[
  \Gamma_2=1+\varepsilon_2.
\]
\end{lemma}

\begin{theorem}[Zero-sum MGF interface]
\label{thm:zero-sum}
\lean{ExchangeableHoeffding.MainProofObligations}
The high-level interface contains the zero-sum statement used by the proof:
if \(a\in\mathbf 1^\perp\), then
\[
  \E \exp(\lambda\langle a,X\rangle)
  \le
  \exp\left\{\frac{\lambda^2}{2}\Gamma_N\|a\|_2^2\right\}.
\]
\end{theorem}

\begin{theorem}[Main exchangeable Hoeffding bound]
\label{thm:main}
\lean{ExchangeableHoeffding.theorem_2_1}
\uses{lem:projection-zero-sum, thm:zero-sum}
The theorem in \texttt{Main.lean} proves the reduction from the zero-sum MGF
interface to the paper's centered weighted statistic.  In Lean form:
\[
\begin{gathered}
\texttt{theorem\_2\_1 (P : MainProofObligations)}\\
\texttt{  (hN : 2 <= N) (hn : n <= N)}
\end{gathered}
\]
returns the MGF bound for all exchangeable \([-1,1]\)-bounded laws, all
weights, and all \(\lambda\in\R\).
\end{theorem}

\begin{corollary}[Tail bound]
\lean{ExchangeableHoeffding.tailThreshold, ExchangeableHoeffding.corollary_2_2}
The corollary uses the threshold
\[
  \|P_{\mathbf 1^\perp}\widetilde w\|_2
  \sqrt{2\Gamma_N\log(1/\delta)}
\]
and is exported from the high-level interface as \texttt{corollary\_2\_2}.
\end{corollary}

\section{Proof Architecture}

The paper proof is organized as a sequence of reductions, and the Lean files
mirror that sequence.

\subsection{Projection Step}

\[
  a=P_{\mathbf 1^\perp}\widetilde w,\qquad
  \sum_i a_i=0,\qquad
  T_w(x)=\langle a,x\rangle.
\]
This is the responsibility of \texttt{Basic.lean} and the theorem
\texttt{sum\_projOnePerp\_eq\_zero}.  The Lean definition of
\texttt{centeredWeightedSum} already uses the projected dot-product form, so
the main theorem can invoke the zero-sum MGF statement directly.

\subsection{Symmetrization}

Exchangeability gives \(X_\pi\stackrel d=X\).  In Lean this is a statement
about \texttt{Measure.map}.  The remaining integral-level convexity and
vertex-reduction arguments are part of the high-level proof interface.

\subsection{Hamming Slices}

\begin{definition}[Hamming-slice API]
\lean{ExchangeableHoeffding.sliceSubsets, ExchangeableHoeffding.sliceAverage,
ExchangeableHoeffding.elemSym, ExchangeableHoeffding.sliceFunctional,
ExchangeableHoeffding.AtMostTwoValues, ExchangeableHoeffding.SphereSection}
The vertex reduction produces uniform averages over \(k\)-subsets
\(S\subseteq [N]\):
\[
  \log \E_{S_k}\exp\left\{\sum_{i\in S_k}y_i\right\}
  \le
  \frac{\Gamma_N}{8}\|y\|_2^2,
  \qquad
  \sum_i y_i=0.
\]
The functional
\[
  F_{N,k}(y)=
  \log\left\{
    \frac{e_k(e^{y_1},\ldots,e^{y_N})}{\binom Nk}
  \right\}
\]
is represented by \texttt{sliceFunctional}.
\end{definition}

\begin{proposition}[Two-level slice extremizer]
\label{prop:two-level}
\lean{ExchangeableHoeffding.SliceProofObligations}
The field \texttt{proposition\_4\_4} states that the maximum of
\(F_{N,k}\) on a zero-sum sphere is attained by a vector with at most two
coordinate values.
\end{proposition}

\begin{lemma}[Slice inequality]
\label{lem:slice}
\lean{ExchangeableHoeffding.MainProofObligations}
\uses{prop:two-level}
The high-level field \texttt{slice\_inequality} records the Hamming-slice MGF
bound needed for the proof of Theorem~\ref{thm:main}.
\end{lemma}

\subsection{Three-Coordinate Reduction}

\begin{definition}[Three-coordinate objects]
\lean{ExchangeableHoeffding.threeObjective,
ExchangeableHoeffding.ThreeSphereSection,
ExchangeableHoeffding.HasDuplicateCoordinate,
ExchangeableHoeffding.ThreePointProofObligations}
The file \texttt{ThreePoint.lean} formalizes the objects behind Lemmas 4.2 and
4.3.  The constrained objective is
\[
  A\sum_{i=1}^{3}e^{x_i}
  +
  B\sum_{i=1}^{3}e^{-x_i},
\]
on the section \(x_1+x_2+x_3=s\),
\(x_1^2+x_2^2+x_3^2=q\).  The interface states both the Lagrange sign lemma
and the conclusion that a maximizer cannot have three distinct coordinates.
\end{definition}

\subsection{Hypergeometric MGF}

\begin{definition}[Hypergeometric objects]
\lean{ExchangeableHoeffding.hypergeomPMF,
ExchangeableHoeffding.hypergeomMgf,
ExchangeableHoeffding.hypergeomB,
ExchangeableHoeffding.HypergeometricProofObligations}
The two-level reduction leaves a centered hypergeometric count.  The file
\texttt{Hypergeometric.lean} defines the PMF, the MGF, and
\[
  B_{N,m}
  =
  (N-s)^2\sum_{\ell=N-s}^{N-1}\ell^{-2},
  \qquad s=m\wedge(N-m).
\]
The interface contains Hoeffding's lemma and the sharpened hypergeometric MGF
bound
\[
  \log \E e^{t(H-mK/N)}
  \le
  \frac{t^2}{8}B_{N,m}.
\]
\end{definition}

\section{Sharpness and Comparison}

\begin{lemma}[Gamma algebra]
\lean{ExchangeableHoeffding.GammaProofObligations,
ExchangeableHoeffding.gamma_lt_barber}
The \texttt{GammaProofObligations} interface records the closed form,
asymptotic expansion, and comparison
\[
  \Gamma_N<1+\varepsilon_N\qquad (N\ge 3).
\]
The exported theorem \texttt{gamma\_lt\_barber} derives this comparison from
that interface.
\end{lemma}

\begin{proposition}[Optimality lower bound]
\lean{ExchangeableHoeffding.AdmissibleConstant,
ExchangeableHoeffding.optimalLowerBound,
ExchangeableHoeffding.optimal_lower_bound}
The lower-bound part of the paper is represented by the admissible-constant
predicate.  Any replacement constant \(C\) valid uniformly in the main theorem
must satisfy
\[
  C \ge
  \begin{cases}
    N/(N-1), & N \text{ even},\\
    (N+1)/N, & N \text{ odd}.
  \end{cases}
\]
\end{proposition}

\section{Reading Guide}

For a first pass through the project, read the Lean files in this order:
\[
\texttt{Basic}
\longrightarrow
\texttt{Gamma}
\longrightarrow
\texttt{Hypergeometric}
\longrightarrow
\texttt{ThreePoint}
\longrightarrow
\texttt{HammingSlice}
\longrightarrow
\texttt{Main}.
\]
This order follows the paper from notation and centering, through the analytic
ingredients, to the final theorem and sharpness statements.
